\section*{Chapter \ref{ch:g4:algorithms-repeat-and-decide} --- Algorithms: Repeat and Decide}

\begin{solution}{exo:g4:algorithms-repeat-and-decide:1}
It walks 2 squares right, 2 squares up, then 2 squares left: three sides of a
square. It stops on the third square up of the leftmost column (2 squares
above the start), facing down.
\end{solution}

\begin{solution}{exo:g4:algorithms-repeat-and-decide:2}
\begin{quote}
\textbf{repeat} 2 times:\\
\hspace*{1.5em}forward, forward, forward, turn left, forward, turn left
\end{quote}
The robot walks along a row of 4 squares, goes up one, comes back along the
row above, and goes down one: it is back on its start square, facing right.
\end{solution}

\begin{solution}{exo:g4:algorithms-repeat-and-decide:3}
\begin{center}
\begin{tabular}{l*{7}{c}}
card looked at & 5 & 11 & 2 & 8 & 14 & 6 & 9 \\
best so far & 5 & 11 & 11 & 11 & 14 & 14 & 14
\end{tabular}
\end{center}
The largest card is 14.
\end{solution}

\begin{solution}{exo:g4:algorithms-repeat-and-decide:4}
Replace ``larger'' by ``smaller'' (and read ``best so far'' as ``smallest so
far''). The smallest card is 2.
\end{solution}

\begin{solution}{exo:g4:algorithms-repeat-and-decide:5}
(a), (c) and (e) are \omterm{def:g4:algorithms-repeat-and-decide:condition}{conditions}: their answer is yes or no. (b) and (f) are
\omterm{def:g1:instructions-in-order:instruction}{instructions}. (d) is a question, but its answer is a number, so it is not a
\omterm{def:g4:algorithms-repeat-and-decide:condition}{condition}.
\end{solution}

\begin{solution}{exo:g4:algorithms-repeat-and-decide:6}
Step 1: the second square is free, forward. Step 2: the square ahead is the
wall, turn left (now facing up). Steps 3 and 4: forward, forward. The robot
stops on the second square of the third row, facing up.
\end{solution}

\begin{solution}{exo:g4:algorithms-repeat-and-decide:7}
Units: $6 + 7 = 13$; write 3 and carry 1 to the tens. Tens: $5 + 2 + 1 = 8$;
write 8. There is no column left: the sum is 83.
\end{solution}

\begin{solution}{exo:g4:algorithms-repeat-and-decide:8}
``Forward, forward, turn left'': after the first two steps it turns up, so it
walks around the square above and to the right of S. With ``turn right'', it
would turn down after two steps and walk around a square below S, which is off
this grid.
\end{solution}

\begin{solution}{exo:g4:algorithms-repeat-and-decide:9}
(a) 5 \omterm{def:g1:instructions-in-order:instruction}{instructions}: 2 right, turn, 2 up. (b) 6 \omterm{def:g1:instructions-in-order:instruction}{instructions}: 2 up, turn,
2 right. (c) 8 \omterm{def:g1:instructions-in-order:instruction}{instructions}: 3 right, 2 up, 1 left. All three reach the goal;
(a) is the shortest.
\end{solution}

\begin{solution}{exo:g4:algorithms-repeat-and-decide:10}
The robot walks along the edge of the grid: right to the corner, up, left, down,
and back to its start square, turning left at each corner. It never passes
through the star, which is inside, so it goes round and round forever. An
\omterm{def:g4:algorithms-repeat-and-decide:algorithm}{algorithm} must always come to an end: this one needs, for example, a way to
reach squares that are not on the edge, or a rule that stops it when it is back
where it started.
\end{solution}

\begin{solution}{exo:g4:algorithms-repeat-and-decide:11}
For example:
\begin{quote}
Start with a count of 0.\\
\textbf{repeat} for each card of the row, from left to right:\\
\hspace*{1.5em}\textbf{if} the card is red \textbf{then} add 1 to the count\\
When there is no card left, the count is the number of red cards.
\end{quote}
\end{solution}
